%@ pages 324-327
%@ status translated
% Introduction (pp. 324-327).

\chapter*{Introduction}
\addcontentsline{toc}{chapter}{Introduction}
\markboth{Abelian Categories}{Introduction}

This work contains some general theorems on abelian categories and some
applications of these theorems to the study of modules. We have devoted a
good deal of space to recalls and to the statement of elementary properties
of abelian categories; there are two reasons for this: the first is that the
subject is recent enough and that we want to give the statements the form
that suits us.

The second reason is that we should like to convince the non-specialists and
to offer them an overall exposition: let us say therefore to this group of
readers that abelian categories were introduced by \textsc{Buchsbaum} and
\textsc{Grothendieck} in order to generalise the homological methods of
\textsc{Cartan} and \textsc{Eilenberg} \cite{cartan-eilenberg}. A comparison will immediately
show the interest of this notion: let $\{M_i\}_{i\in I}$ be a family of
modules over a ring $A$, let $T_i$ be the lattice of submodules of $M_i$
($i\in I$) and $\operatorname{Hom}_A(M_i,M_j)$ the abelian group formed by
the $A$-linear maps of $M_i$ into $M_j$ ($i,j\in I$). One of the aims of
lattice theory is to study the family of the modules $M_i$ by means of the
family of the lattices $T_i$ \cite{chevalley-grothendieck}; likewise, the aim of homological algebra is
the study of the $M_i$ by means of the groups
$\operatorname{Hom}_A(M_i,M_j)$ and of the laws of composition which connect
these groups. For a suitable choice of the family $(M_i)_{i\in I}$, the
datum of the $\operatorname{Hom}_A(M_i,M_j)$ together with the laws of
composition permits one to reconstruct the lattices $T_i$ ($T_i$ is the
lattice of the ``subobjects'' of $M_i$). The data are therefore richer in
homological algebra; it follows that the results are more precise.

The results to which we are led apply principally to the following special
cases:

\begin{description}
\item[\normalfont a.] \emph{Category of right modules over a right noetherian
ring.} --- Our statements are then translated into the classical language of
module theory; we make this translation in part in order to conform to
established traditions. There is nevertheless no doubt that the study of
module categories brings more general abelian categories into play (see the
notion of quotient category).

\item[\normalfont b.] \emph{Category of quasi-coherent sheaves on a noetherian
prescheme} (this application is treated in Chapter VI).

\item[\normalfont c.] \emph{Category of commutative algebraic groups} (see
\cite{serre-groupes}; this application will be the subject of a later publication).

\item[\normalfont d.] \emph{Category of Hopf algebras over a field,
connected, with associative and commutative multiplication and diagonal}
(this application does not deserve a publication).
\end{description}

Chapter~I consists of recalls and supplements to the literature published up
to the present day. We insist especially on the notion of \emph{equivalence}
of two categories: two categories $\mathbf{A}$ and $\mathbf{B}$ are
equivalent if they are isomorphic when the morphisms are defined as follows: a
morphism of $\mathbf{A}$ into $\mathbf{B}$ is the class of the functors of
$\mathbf{A}$ into $\mathbf{B}$ which are isomorphic to a given functor. If
one adopts this point of view, one obtains, for example, the following result:
let $A$ be a ring with identity such that every unitary right $A$-module which
is projective and of finite type is free; let $\mathbf{A}$ (resp.\
$\mathbf{B}$) be the category of unitary right modules over $A$ [resp.\ over
the ring $\mathbf{M}_n(A)$ formed of the $(n,n)$ matrices with coefficients in
$A$]; the group of automorphisms of $\mathbf{A}$ (resp.\ of $\mathbf{B}$) is
the quotient of the group of automorphisms of the ring $A$ [resp.\ of the
ring $\mathbf{M}_n(A)$] by the subgroup of inner automorphisms. Since the
categories $\mathbf{A}$ and $\mathbf{B}$ are equivalent (cf.\ Chapter~V,
\S1), we recover that the quotients in question are isomorphic.

Chapter~II presents the material which we work with in what follows: let us
say, roughly, that an abelian category is \emph{noetherian} if all the objects
of that category are noetherian, that an abelian category is \emph{locally
noetherian} if there are inductive limits, if the inductive limit functor is
exact and if every object is an inductive limit of noetherian objects. We
show that every noetherian category $\mathbf{A}$ is equivalent to the category
of noetherian objects of a locally noetherian category $\mathbf{B}$.
Conversely, $\mathbf{B}$ is equivalent to the category of additive,
contravariant, left exact functors of $\mathbf{A}$ into the category of abelian
groups. The category $\mathbf{B}$ is therefore equivalent to a quotient
category of the category of all additive contravariant functors of
$\mathbf{A}$ into abelian groups; it follows that $\mathbf{B}$ is
``almost'' a category of modules (cf.\ Chapter~II, \S1).

In Chapter~III we are interested in the ``language modulo $\mathcal{C}$'' of
\textsc{Serre}: let $\mathbf{A}$ and $\mathbf{B}$ be two abelian categories,
$T$ and $S$ two additive functors, $T:\mathbf{A}\to\mathbf{B}$,
$S:\mathbf{B}\to\mathbf{A}$. We suppose that $T$ is exact, that $S$ is
adjoint to $T$ and that $T\circ S$ is isomorphic to the identity functor of
$\mathbf{B}$. Under these conditions the category $\operatorname{Ker} T$,
which is formed of the objects $M$ of $\mathbf{A}$ such that $TM$ is zero, is
thick \cite{grothendieck-1}; moreover, $T$, defined by passing to the quotient, induces an
equivalence between the quotient category of $\mathbf{A}$ by
$\operatorname{Ker} T$ \cite{grothendieck-1} and the category $\mathbf{B}$. Conversely, let
$\mathbf{C}$ be a thick subcategory of $\mathbf{A}$ and $T$ the canonical
functor from $\mathbf{A}$ into $\mathbf{A}/\mathbf{C}$ \cite{grothendieck-1}. We say that
$\mathbf{C}$ is a \emph{localizing subcategory} of $\mathbf{A}$ if there
exists a functor adjoint to $T$.

When $\mathbf{A}$ is the category of modules over a commutative ring $A$, we
can give the following example: let $\Sigma$ be a multiplicatively closed subset
of $A$, $\mathbf{B}$ the category of modules over the ring $A_\Sigma$ and $T$
the functor which associates to every $A$-module $M$ its ``localised''
$M_\Sigma$. The functor $T$, defined by passing to the quotient, induces an
equivalence between $\mathbf{A}/\operatorname{Ker} T$ and $\mathbf{B}$.

Chapter~IV contains a few results on locally noetherian categories: we prove
that an injective object of such a category $\mathbf{A}$ is a direct sum of a
family of indecomposable injective objects. We also deal with the Krull
dimension of $\mathbf{A}$; for this we ``filter'' $\mathbf{A}$ by means of
localizing subcategories $\mathbf{A}_0\subset\mathbf{A}_1\subset\mathbf{A}_2\subset\cdots$
such that the quotients $\mathbf{A}_{i+1}/\mathbf{A}_i$ are locally finite
categories (that is to say, $\mathbf{A}_{i+1}/\mathbf{A}_i$ is locally
noetherian and every object is an inductive limit of objects of finite
length); we associate to these locally finite categories complete topological
rings, playing the same role as complete local rings do in commutative
algebra.

Chapter~V deals with applications to module theory. When $\mathbf{A}$ contains
in its centre a commutative noetherian ring $R$ and $M$ is a finitely
generated $R$-module, we make explicit the constructions proposed in
Chapter~IV. We show, in particular, that the study of the categories
$\mathbf{A}_{i+1}/\mathbf{A}_i$ is equivalent to the study of the completions
of $A$ for certain topologies. This chapter sets out a plan for attacking
non-commutative noetherian rings. It remains to be seen whether the methods
envisaged are practicable in the study of usual rings. In this connection let
us point out a ring which \textsc{Dieudonn\'e} has introduced into the study of
formal groups:

\begin{quote}
Let $k$ be an algebraically closed field of characteristic $p>0$, $W$ the ring
of Witt vectors with coefficients in $k$ and $\sigma$ the Frobenius
automorphism of $W$. We denote by $A$ the ring whose elements are the formal
series of the form
\[
  a = w + \sum_{r\geq 1} a_r\, F^{r} + \sum_{s\geq 1} b_s\, V^{s},
\]
where $w$, $a_r$ and $b_s$ range over the ring $W$; multiplication is defined
by the following rules:
\[
  V \cdot F = F \cdot V = p, \qquad
  F \cdot w = \sigma(w) \cdot F \qquad\text{and}\qquad
  w \cdot V = V \cdot \sigma(w)
  \quad\text{if } w \in W.
\]
\end{quote}

The ring $A$ is noetherian (on the right and on the left). If we denote by
$\mathbf{A}$ the category of right $A$-modules, then:

\begin{itemize}
\item Every simple object of $\mathbf{A}$ is isomorphic to
      $A/(F \cdot A + V \cdot A)$.
\item If $\mathbf{A}_0$ is the smallest localizing subcategory of
      $\mathbf{A}$ which contains $A/(F \cdot A + V \cdot A)$, then every
      simple object of $\mathbf{A}/\mathbf{A}_0$ is of one of the following
      types: $A/F \cdot A$, $A/V \cdot A$ or $A/\big((V^{r} - F^{q})\cdot
      A\big)$ with $(r,q) = 1$.
\item If $\mathbf{A}_1$ is the smallest localizing subcategory of
      $\mathbf{A}$ which contains the above modules, then every simple object
      of $\mathbf{A}/\mathbf{A}_1$ is isomorphic to $A$.
\end{itemize}

This work has already been presented in three lectures:

\begin{itemize}
\item \emph{Injective objects in abelian categories},
      \emph{S\'eminaire Dubreil--Pisot}, t.~12, 1958--1959, n\textsuperscript{o}~17,
      32 pages.
\item \emph{Localisation in non-commutative rings}, \emph{S\'eminaire
      Dubreil--Pisot}, t.~13, 1959--1960, n\textsuperscript{o}~2, 35 pages.
\item \emph{On locally noetherian categories and their applications to the
      algebras studied by Dieudonn\'e} (formal groups), \emph{S\'eminaire
      J.-P. Serre}, 1959--1960.
\end{itemize}